% =============================================================================
% CompetitiveExam/CS401/12-pipelining-risc-cisc-questions.tex
% Competitive Exam Practice 12 - Instruction Pipelining and RISC vs. CISC
%
% No real PYQ is included: the local Week 12 scan contains relevant headings,
% but OCR did not recover a complete stem and answer independently. Per the
% sourcing protocol, those candidates are not presented as authentic PYQs.
%
% Compile from CompetitiveExam/CS401 on Windows/MiKTeX:
%   $env:TEXINPUTS = "../../Preambles;"
%   $env:BIBINPUTS = "../..;"
%   xelatex -interaction=nonstopmode 12-pipelining-risc-cisc-questions.tex
% =============================================================================

\documentclass[11pt]{article}
\usepackage[margin=1in]{geometry}
% =============================================================================
% Preambles/header.tex — shared LaTeX/Beamer preamble for this project
%
% Use from a Beamer lecture by prepending:
%     \input{header}
% and compiling with TEXINPUTS=../Preambles:$TEXINPUTS (the /compile-latex
% skill does this automatically).
%
% PALETTE CONTRACT. Color names defined here MUST match the names in
% Quarto/theme-template.scss so Beamer and Quarto renderings use the same
% palette. The scripts/check-palette-sync.sh script enforces this.
%
% To customize: edit the HEX values below AND the matching $variable in
% Quarto/theme-template.scss, then run scripts/check-palette-sync.sh.
% =============================================================================

% -----------------------------------------------------------------------------
% Engine detection + encoding
%
% Our /compile-latex skill uses XeLaTeX, where inputenc is unnecessary and
% can produce encoding warnings. Only load inputenc under pdfTeX.
% -----------------------------------------------------------------------------
\usepackage{iftex}
\ifPDFTeX
  \usepackage[utf8]{inputenc}
\fi

\usepackage{amsmath, amssymb, amsthm}
\usepackage{booktabs}
\usepackage{graphicx}

% -----------------------------------------------------------------------------
% PALETTE — mirrors Quarto/theme-template.scss variable names.
% Load xcolor and define colors BEFORE anything that references them
% (notably hyperref, hypersetup, and the Beamer color assignments).
% -----------------------------------------------------------------------------
\usepackage{xcolor}

\definecolor{primary-blue}{HTML}{012169}      % headings, accents
\definecolor{primary-gold}{HTML}{B9975B}      % emphasis, borders
\definecolor{highlight-yellow}{HTML}{F2A900}  % markers, alerts
\definecolor{light-bg}{HTML}{E8EDF5}          % title / section backgrounds
\definecolor{jet}{HTML}{1A1A1A}               % body text

% Semantic colors (used in diagrams and annotations)
\definecolor{positive}{HTML}{15803D}          % good / observed / identified
\definecolor{negative}{HTML}{B91C1C}          % bad / problematic / bias
\definecolor{neutral}{HTML}{525252}           % reference / context

% Highlight accents (match .hi-* classes in the SCSS)
\definecolor{hi-slate}{HTML}{314F4F}
\definecolor{hi-green}{HTML}{2E8B57}
\definecolor{hi-red}{HTML}{C41E3A}

% Semantic aliases so skill templates and snippets can write
% \textcolor{accent}{...} without hard-coding "primary-blue".
\colorlet{accent}{primary-blue}
\colorlet{accent2}{primary-gold}

% -----------------------------------------------------------------------------
% Hyperref — loaded AFTER xcolor + palette so link colors resolve.
% -----------------------------------------------------------------------------
\usepackage{hyperref}
\hypersetup{colorlinks=true, linkcolor=primary-blue, urlcolor=primary-blue,
            citecolor=primary-blue, pdfborder={0 0 0}}

% -----------------------------------------------------------------------------
% Course code — set once per deck (after \input{header}, before \begin{document})
% via \coursecode{CS401}. Shown in the footer on every slide so a deck's course
% is identifiable without opening the title page. Empty by default.
% -----------------------------------------------------------------------------
\makeatletter
\newcommand{\coursecode}[1]{\gdef\@coursecodetext{#1}}
\gdef\@coursecodetext{}
\makeatother

% -----------------------------------------------------------------------------
% Beamer theme assignments (only apply under Beamer).
%
% We detect Beamer via \@ifclassloaded{beamer}, which is the documented,
% non-internal way. This must be wrapped in \makeatletter/\makeatother
% because \@ifclassloaded contains an @-catcode character.
% -----------------------------------------------------------------------------
\makeatletter
\@ifclassloaded{beamer}{%
  \usefonttheme{professionalfonts}
  \setbeamercolor{structure}{fg=primary-blue}
  \setbeamercolor{frametitle}{fg=primary-blue}
  \setbeamercolor{title}{fg=primary-blue}
  \setbeamercolor{subtitle}{fg=primary-gold}
  \setbeamercolor{author}{fg=primary-blue}
  \setbeamercolor{institute}{fg=neutral}
  \setbeamercolor{date}{fg=primary-gold}
  \setbeamercolor{itemize item}{fg=primary-blue}
  \setbeamercolor{itemize subitem}{fg=primary-gold}
  \setbeamercolor{alerted text}{fg=primary-gold}
  \setbeamercolor{block title}{fg=white, bg=primary-blue}
  \setbeamercolor{block body}{bg=light-bg}
  % Semantic block variants: block=definition/concept (blue), exampleblock=
  % worked example (green/positive), alertblock=key takeaway or caution (gold).
  % Keep to <=2 colored boxes per slide across ALL three variants (INV-7).
  \setbeamercolor{block title example}{fg=white, bg=positive}
  \setbeamercolor{block body example}{bg=light-bg}
  \setbeamercolor{block title alerted}{fg=white, bg=primary-gold}
  \setbeamercolor{block body alerted}{bg=light-bg}

  % Minimal footer: course code (left) + page X / N (right), no navigation clutter
  \setbeamertemplate{navigation symbols}{}
  \setbeamertemplate{footline}{%
    \color{neutral}\footnotesize\hspace{0.7em}\@coursecodetext\hfill
    \insertframenumber/\inserttotalframenumber\hspace{0.7em}\vspace{0.3em}}
}{}
\makeatother

% -----------------------------------------------------------------------------
% TikZ setup — libraries the snippet gallery uses
% -----------------------------------------------------------------------------
\usepackage{tikz}
\usetikzlibrary{arrows.meta, positioning, calc, decorations.pathreplacing,
                fit, shapes.geometric, backgrounds}

% Shared TikZ styles — snippets and hand-written diagrams can pick these up
% via `\begin{tikzpicture}[every node/.style={...}]` or by naming specific
% styles (e.g., `\node[dag-node] (X) at (0,0) {X};`).
\tikzset{
  % Node styles for causal DAGs and similar
  dag-node/.style = {
    draw=primary-blue, fill=light-bg, rounded corners=2pt,
    minimum width=1.2cm, minimum height=0.9cm, align=center, font=\small
  },
  decision-node/.style = {
    draw=primary-gold, fill=white, diamond, aspect=1.8,
    minimum width=1.8cm, minimum height=1.2cm, align=center, font=\small,
    inner sep=1pt
  },
  % Edge styles — observed vs counterfactual
  observed-edge/.style = {-{Latex[length=2.5mm]}, thick, draw=jet},
  counterfactual-edge/.style = {-{Latex[length=2.5mm]}, thick, draw=neutral, dashed},
  confound-edge/.style = {-{Latex[length=2.5mm]}, thick, draw=negative, dashed},
  % Data-point styles for plots
  observed-dot/.style = {fill=primary-blue, draw=primary-blue, circle, inner sep=1.2pt},
  counterfactual-dot/.style = {fill=white, draw=neutral, circle, inner sep=1.2pt}
}

% -----------------------------------------------------------------------------
% Convenience macros
% -----------------------------------------------------------------------------
\newcommand{\muted}[1]{\textcolor{neutral}{#1}}
\newcommand{\key}[1]{\textcolor{primary-gold}{\textbf{#1}}}
\newcommand{\good}[1]{\textcolor{positive}{#1}}
\newcommand{\bad}[1]{\textcolor{negative}{#1}}

% Standout frame macro: full-bleed dark background for section transitions.
% Beamer-only (uses \begin{frame}/\setbeamercolor/\usebeamerfont) -- do not
% call from an article-class file (e.g. Notes/<CODE>/*-notes.tex). \muted,
% \key, \good, \bad, and \coursecode above are plain xcolor/gdef and are
% safe to use from either class; verified when Notes/article-class support
% was added.
% Expand: \transitionslide{Your transition title}
\newcommand{\transitionslide}[1]{%
  \begingroup
  \setbeamercolor{background canvas}{bg=primary-blue}
  \begin{frame}[plain]
    \centering
    \vfill
    {\usebeamerfont{title}\color{white}\Large #1\par}
    \vfill
  \end{frame}
  \endgroup
}

% Section divider frame (Beamer-only), an alternative to \transitionslide with a
% darker two-line style: near-black (jet) background, a small white label line
% above a larger gold title, and a thin gold rule along the bottom edge.
% Expand: \sectiondivider{Section 2}{Pointers}
\newcommand{\sectiondivider}[2]{%
  \begingroup
  \setbeamercolor{background canvas}{bg=jet}
  \begin{frame}[plain]
    \vspace*{3.0cm}
    {\color{white}\ttfamily\large #1\par}
    \vspace{0.5cm}
    {\color{primary-gold}\LARGE #2\par}
    \begin{tikzpicture}[remember picture, overlay]
      \draw[primary-gold, line width=2.5pt]
        ($(current page.south west)+(0,0.1)$) -- ($(current page.south east)+(0,0.1)$);
    \end{tikzpicture}
  \end{frame}
  \endgroup
}

% -----------------------------------------------------------------------------
% Channel watermark — "Concepts That Click" (YouTube).
%
% Article-class files (CompetitiveExam Q/A sets, Notes, Assignments) get a
% light diagonal draft watermark on every page plus a centered footer naming
% the channel. Beamer decks get a subtle TikZ background watermark instead
% (the footline already carries the course code). Centralized here so every
% MCQ PDF is branded without per-file edits.
% -----------------------------------------------------------------------------
\makeatletter
\@ifclassloaded{article}{%
  \usepackage{draftwatermark}
  \SetWatermarkText{Concepts That Click}
  \SetWatermarkScale{1}
  \SetWatermarkFontSize{2cm}
  \SetWatermarkLightness{0.86}
  \SetWatermarkAngle{45}
  \usepackage{fancyhdr}
  \pagestyle{fancy}
  \fancyhf{}
  \fancyfoot[C]{\footnotesize\textcolor{neutral}{Concepts That Click~|~YouTube: \href{https://www.youtube.com/@concepts.thatclick}{@concepts.thatclick}}}
  \renewcommand{\headrulewidth}{0pt}
  \renewcommand{\footrulewidth}{0pt}
  \fancypagestyle{plain}{%
    \fancyhf{}%
    \fancyfoot[C]{\footnotesize\textcolor{neutral}{Concepts That Click~|~YouTube: \href{https://www.youtube.com/@concepts.thatclick}{@concepts.thatclick}}}%
    \renewcommand{\headrulewidth}{0pt}%
    \renewcommand{\footrulewidth}{0pt}%
  }
}{}
\@ifclassloaded{beamer}{%
  \setbeamertemplate{background}{%
    \begin{tikzpicture}[remember picture, overlay]
      \node[rotate=45, opacity=0.055, align=center]
        at (current page.center)
        {\Huge\textcolor{neutral}{Concepts That Click}};
    \end{tikzpicture}%
  }
}{}
\makeatother 
\coursecode{CS401}
\renewcommand{\thesection}{12.\arabic{section}}

\title{Competitive Exam Practice 12: Instruction Pipelining and RISC vs. CISC\\
  \large GATE (Computer Science) Pattern}
\author{Auqib Hamid Lone\\[0.3em]
  \emph{Department of Computer Science \& Engineering}, \emph{GCET Ganderbal}}
\date{\today}

\begin{document}
\maketitle

\noindent Each question carries 1--2 marks. MCQ = single correct option;
MSQ = multiple-select; NAT = Numerical Answer Type. Every item below is
\textit{[Original, GATE-pattern]}. No item is labelled as a real PYQ because
the local scanned PYQ pages did not yield a complete, independently verifiable
stem and answer.

\begin{enumerate}
\item[\textbf{Q1}] \textit{[Original, GATE-pattern]} \textbf{[NAT, 2 marks]}
Twelve independent instructions execute in an ideal balanced five-stage
pipeline with one cycle per stage and no stalls. What are the pipelined
completion time and the finite-stream speedup over a non-pipelined execution,
respectively? Report the speedup to two decimal places.

\item[\textbf{Q2}] \textit{[Original, GATE-pattern]} \textbf{[NAT, 2 marks]}
Consider the five-stage pipeline \texttt{IF, ID, EX, MEM, WB} and the sequence
\texttt{LW R1,0(R2)}, \texttt{ADD R3,R1,R4}, and \texttt{SUB R5,R3,R6}.
Load data is available after MEM; ALU results are forwarded at the end of EX.
Starting the first IF in cycle 1, give the cycle in which \texttt{SUB} completes
WB and the number of bubbles required, in the form ``cycle, bubbles''.

\item[\textbf{Q3}] \textit{[Original, GATE-pattern]} \textbf{[MCQ, 1 mark]}
An ideal five-stage pipeline is full and has no hazards. Which statement is
correct for a long independent instruction stream?
\begin{enumerate}
  \item[(A)] Each instruction has one-cycle latency and one-cycle throughput.
  \item[(B)] Each instruction has approximately five-cycle latency, while
    the stream can complete one instruction per cycle.
  \item[(C)] Both latency and throughput are approximately five cycles per
    instruction.
  \item[(D)] Throughput improves only if the pipeline has one stage.
\end{enumerate}

\item[\textbf{Q4}] \textit{[Original, GATE-pattern]} \textbf{[NAT, 1 mark]}
Branches are $20\%$ of the dynamic instruction stream. The predictor
mispredicts $10\%$ of branches, and each misprediction costs 3 cycles. Ignoring
all other effects, what is the branch contribution to CPI? Report to two
decimal places.

\item[\textbf{Q5}] \textit{[Original, GATE-pattern]} \textbf{[MCQ, 1 mark]}
In a five-stage pipeline, an instruction in MEM and a younger instruction in
IF request the same single-ported unified memory in the same cycle. Which pair
is the most direct classification and one valid remedy?
\begin{enumerate}
  \item[(A)] Data hazard; forward the producer's ALU result
  \item[(B)] Structural hazard; split instruction and data memories
  \item[(C)] Control hazard; rename the destination register
  \item[(D)] Data hazard; add a branch-prediction entry
\end{enumerate}

\begin{sloppypar}
\item[\textbf{Q6}] \textit{[Original, GATE-pattern]} \textbf{[MSQ, 1 mark]}
Which statements are correct about the RISC/CISC comparison in a fixed
workload?
\begin{enumerate}
  \item[(A)] RISC often favors regular formats and register-centered
    execution.
  \item[(B)] CISC can express more work per architectural instruction.
  \item[(C)] RISC always has fewer instructions and CISC always has lower CPI.
  \item[(D)] Total time should be compared using instruction count, CPI, and
    clock rate, rather than an ISA label alone.
\end{enumerate}
\end{sloppypar}

\item[\textbf{Q7}] \textit{[Original, GATE-pattern]} \textbf{[MCQ, 1 mark]}
For the same workload, a RISC-like implementation executes
$1.20\times10^9$ instructions at CPI $=1.1$ and $f=2.4$~GHz. A CISC-like
implementation executes $0.80\times10^9$ instructions at CPI $=2.0$ and
$f=3.0$~GHz. Which statement is correct?
\begin{enumerate}
  \item[(A)] RISC-like is faster because its CPI is lower.
  \item[(B)] CISC-like is faster because its total time is approximately
    $0.533$~s versus $0.550$~s.
  \item[(C)] Both take the same time because CPI is the only relevant metric.
  \item[(D)] No comparison is possible without knowing the instruction names.
\end{enumerate}

\item[\textbf{Q8}] \textit{[Original, GATE-pattern]} \textbf{[MSQ, 1 mark]}
Which statements correctly identify parallelism or workload constraints?
\begin{enumerate}
  \item[(A)] One instruction applied to eight data elements is SIMD-style
    data parallelism.
  \item[(B)] Independent cores executing different tiles are MIMD-style
    parallelism.
  \item[(C)] A large working set can make cache and virtual-memory behavior
    relevant even when the front end is pipelined.
  \item[(D)] A higher clock rate alone proves that an implementation is
    faster for every workload.
\end{enumerate}
\end{enumerate}

\end{document}